\chapter*{Introducción} % sin número
\addcontentsline{toc}{chapter}{Introducción}

El estudio del clima ha sido una preocupación constante a lo largo de la historia de la humanidad. Desde las antiguas civilizaciones, que basaban sus predicciones en la observación empírica de fenómenos naturales —como la formación de nubes, los patrones de viento o los cambios de temperatura \citep{ossendrijver2021babylonian}—, hasta los científicos contemporáneos que emplean sofisticados modelos físico-matemáticos para describir el comportamiento atmosférico \citep{lorenz1963deterministic}, la necesidad de anticipar las condiciones climáticas ha persistido como un desafío fundamental.

En la actualidad, la capacidad de predecir variables meteorológicas con precisión es vital para la planificación urbana, la agricultura, la gestión de recursos hídricos, la aviación, la generación de energía eólica y la prevención de desastres naturales, como tormentas o inundaciones \citep{wmo2023state}. No obstante, la atmósfera es un sistema caótico gobernado por ecuaciones no lineales, lo que ocasiona que pequeñas variaciones en las condiciones iniciales generen divergencias significativas en las predicciones; este fenómeno es conocido como el “efecto mariposa” \citep{lorenz1972predictability}. Esta complejidad inherente convierte la reconstrucción y proyección de estados climáticos en una tarea computacional y teóricamente exigente.

Con el desarrollo de sistemas de computación de alto rendimiento, la recolección de datos meteorológicos ha experimentado una revolución. Estaciones terrestres, boyas oceánicas, radares Doppler y satélites de observación —como los programas ERA5 del ECMWF o GOES de la NASA/NOAA \citep{aragao2022cyclonic}— proporcionan flujos masivos de información en tiempo real, capturando variables como temperatura, presión atmosférica, velocidad y dirección del viento, humedad y radiación solar \citep{hasler1989automatic}.

Sin embargo, los métodos tradicionales para procesar estos datos, basados en simulaciones numéricas como los Modelos de Circulación General (GCMs) \citep{guilyardi2009coupled} o el Weather Research and Forecasting Model (WRF) \citep{skamarock2001wrf}, enfrentan dos limitaciones principales:
\begin{enumerate}
    \item Alto costo computacional: La necesidad de resolver ecuaciones diferenciales complejas, como las de Navier-Stokes, en una malla tridimensional con alta resolución espacial y temporal.

    \item Dependencia de parametrizaciones empíricas: Muchos procesos físicos, como la formación de nubes o la turbulencia, ocurren a escalas menores que la resolución de la malla del modelo. Esto obliga a realizar aproximaciones mediante parametrizaciones basadas en relaciones simplificadas, lo que introduce incertidumbre en los resultados.
\end{enumerate}

En este escenario destaca el potencial de las Redes Neuronales Físicamente Informadas (PINNs), una rama de la inteligencia artificial que busca cerrar la brecha entre los modelos basados exclusivamente en datos y las leyes físicas fundamentales. Las PINNs son arquitecturas que incorporan ecuaciones diferenciales parciales (EDPs), como las de Navier-Stokes para fluidos o la de conservación de la energía, directamente en su función de pérdida durante el entrenamiento \citep{cai2021physics}.

A diferencia de las redes neuronales tradicionales, las PINNs no solo aprenden correlaciones estadísticas, sino que garantizan la consistencia con las leyes físicas. Esto permite realizar entrenamientos sin necesidad de volúmenes masivos de datos, ya que el conocimiento físico actúa como un regularizador. Además, resultan sumamente eficaces para resolver problemas inversos, permitiendo inferir parámetros ocultos a partir de observaciones dispersas.

Aunque se ha probado su eficacia en la predicción de flujos turbulentos y la asimilación de datos heterogéneos, su aplicación a escalas locales para la reconstrucción de campos de viento y presión a partir de estaciones meteorológicas dispersas sigue siendo un área poco explorada. Moreno Soto et al. \cite{morenosoto2024pinns} presentaron una de las primeras aplicaciones de PINNs a la reconstrucción meteorológica de alta resolución a partir de estaciones dispersas, sirviendo como referencia metodológica para este trabajo. En este proyecto, el caso de estudio corresponde a la región Caribe de Colombia, a partir de registros de estaciones del IDEAM disponibles en el portal de Datos Abiertos Colombia. Los datos cubren el mes de diciembre de 2025 e incluyen posición geográfica, tiempo, velocidad y dirección del viento, presión y temperatura. A partir de estas observaciones, se plantea una PINN que recibe como entrada dos coordenadas espaciales y una variable temporal $(x, y, t)$, y predice la presión y las componentes cartesianas del viento $(vel_x, vel_y)$ sobre una malla espacio-temporal, con el fin de estimar condiciones atmosféricas en puntos no instrumentados.